\documentclass[14pt,twoside]{article}

% ==================== Paquetes ====================
\usepackage[utf8]{inputenc}
\usepackage[spanish]{babel}
\usepackage{amsmath, amssymb}
\usepackage{graphicx}
\usepackage{geometry}
\usepackage{fancyhdr}
\usepackage[hidelinks]{hyperref}
\usepackage{enumitem}
\usepackage{mathrsfs}
\usepackage{float}
\usepackage{setspace}
\usepackage{parskip}
\usepackage{tikz}
\usepackage{lmodern}
\usepackage{datetime2}
\usepackage{bm}
\usepackage{booktabs}
\usepackage{tabularx}
\usepackage{array}
\usepackage{xcolor}
\usepackage{listings}

% ==================== Estilo de código ====================
\definecolor{pgfondo}{HTML}{F6F8FA}
\definecolor{pgborde}{HTML}{D0D7DE}
\definecolor{pgazul}{HTML}{0550AE}
\definecolor{pgverde}{HTML}{116329}
\definecolor{pggris}{HTML}{57606A}
\definecolor{pgmorado}{HTML}{8250DF}

\lstdefinestyle{pgbase}{
    basicstyle=\ttfamily\footnotesize,
    backgroundcolor=\color{pgfondo},
    frame=single,
    rulecolor=\color{pgborde},
    framesep=6pt,
    breaklines=true,
    breakatwhitespace=false,
    columns=fullflexible,
    keepspaces=true,
    showstringspaces=false,
    tabsize=4,
    captionpos=b,
    aboveskip=1em,
    belowskip=1em,
    literate=
        {á}{{\'a}}1 {é}{{\'e}}1 {í}{{\'i}}1
        {ó}{{\'o}}1 {ú}{{\'u}}1 {ñ}{{\~n}}1
        {Á}{{\'A}}1 {É}{{\'E}}1 {Í}{{\'I}}1
        {Ó}{{\'O}}1 {Ú}{{\'U}}1 {Ñ}{{\~N}}1
        {¿}{{?`}}1 {¡}{{!`}}1
}

\lstdefinestyle{pgpython}{
    style=pgbase,
    language=Python,
    keywordstyle=\color{pgazul}\bfseries,
    commentstyle=\color{pggris}\itshape,
    stringstyle=\color{pgverde},
    emph={np,pd,sklearn,Pipeline,StandardScaler,GridSearchCV,SVC,KNeighborsClassifier,DecisionTreeClassifier,PolynomialFeatures,LinearRegression},
    emphstyle=\color{pgmorado},
    numbers=left,
    numberstyle=\ttfamily\scriptsize\color{pggris},
    numbersep=10pt
}

\lstdefinestyle{pgterminal}{
    style=pgbase,
    numbers=none
}

% ==================== Márgenes y formato ====================
\geometry{
  top=2.5cm,
  bottom=2.5cm,
  left=2cm,
  right=2cm
}
\singlespacing
\renewcommand{\arraystretch}{1.5}

% ==================== Datos del documento ====================
\newcommand{\documenttitle}{Inteligencia artificial generativa}
\newcommand{\materia}{Inteligencia Artificial y Mini-Robots (2023251)}
\newcommand{\mydate}{\DTMdate{2026-10-04}}

% ==================== Encabezado y pie ====================
\pagestyle{fancy}
\fancyhf{}
\fancyhead[RO]{\textbf{\materia}}
\fancyhead[LO]{\textit{\documenttitle}}
\fancyhead[RE]{\textit{\documenttitle}}
\fancyhead[LE]{\textbf{\materia}}
\fancyfoot[LE,RO]{\textit{Universidad Nacional de Colombia}}
\fancyfoot[RE,LO]{\textbf{\thepage}}
\renewcommand{\headrulewidth}{0.4pt}
\renewcommand{\footrulewidth}{0.4pt}

\begin{document}

% ==================== Cabecera principal ====================
\noindent
\begin{minipage}[c]{0.20\textwidth}
    \centering
    \includegraphics[width=0.90\linewidth]{0.img/UN.png}
\end{minipage}\hfill
\begin{minipage}[c]{0.75\textwidth}
    \textbf{\LARGE{\documenttitle}} \\[0.1cm]
    \hrule
    \vspace{0.1cm}
    \textbf{\materia} \\[0.3cm]
    \large{Diego Fernando Bustos Horta, \textit{dbustosh@unal.edu.co}} \\[0.1cm]
    \large{Julián Mateo Vargas Riveros, \textit{jvargasri@unal.edu.co}} \\[0.1cm]
    $\boxed{\text{\mydate}}$
\end{minipage}

\vspace{0.5cm}

% =====================================================================
Repositorio de Github: https://github.com/dokiun/chatbot.git

\section*{Informe de desarrollo --- ejercicios de Inteligencia Artificial}

Este repositorio reúne la solución de los cinco ejercicios propuestos en la sección 7.11 de \texttt{IA.md}. Cada punto tiene sus fuentes, código y evidencia en una carpeta independiente.

\section{Skills para apoyar el trabajo con IA}

Se definieron cuatro skills locales, cada una con un archivo \texttt{SKILL.md} que describe cuándo usarla, sus pasos y su resultado esperado:

\begin{center}
\begin{tabular}{p{0.36\textwidth} p{0.52\textwidth}}
\textbf{Skill} & \textbf{Función} \\
\hline
\texttt{capturar-requerimientos} & Extraer requisitos y mantener trazabilidad con los documentos fuente. \\
\texttt{documentar-ticket} & Convertir una necesidad en un ticket con alcance y criterios de aceptación. \\
\texttt{implementar-y-probar} & Guiar la implementación y su validación. \\
\texttt{explicar-trabajo} & Establecer la estructura de una explicación HTML del trabajo realizado. \\
\end{tabular}
\end{center}

Se aplicaron al documento \texttt{P1.pdf} como ejemplo de captura y organización del trabajo. El resultado está en \texttt{01/P1-Skills.md}, con los aportes separados por skill. En ese ejemplo no se generó una página HTML ni se implementó software: el archivo documenta el plan y la explicación textual.

\section{Chatbot RAG industrial}

Se construyó un chatbot para consultar manuales de \textbf{pulidora, taladro y compresor}. Hay cuatro PDF en \texttt{02/data}, incluidos dos documentos de compresor. Las versiones Markdown procesadas están en \texttt{02/rag/data}.

El flujo toma el texto del manual, lo limpia y divide en fragmentos; luego genera embeddings con \texttt{embeddinggemma} mediante Ollama y los guarda en SQLite. Para responder, combina similitud semántica con coincidencia de términos, filtra por herramienta cuando la pregunta la identifica y envía los fragmentos recuperados a \texttt{gemma3:1b}. La respuesta incluye el nombre del documento y el número de fragmento.

\textbf{Código y evidencia:}

\begin{itemize}
\item Pipeline, consulta y uso: \texttt{02/rag/README.md}.
\item Preguntas y respuestas generadas: \texttt{02/preguntas\_y\_respuestas.md}.
\end{itemize}

Ejemplo de ejecución, desde la raíz del repositorio:

\begin{verbatim}
python3 02/rag/scripts/build_index.py
python3 02/rag/scripts/embed_with_ollama.py embeddinggemma 2
python3 02/rag/scripts/chat.py "¿Qué debo revisar antes de encender la pulidora?"
\end{verbatim}

La calidad de las respuestas aún requiere revisión humana: algunos manuales extraídos contienen texto desordenado o varios idiomas, y la evaluación registra casos con fuentes mezcladas o respuestas incompletas. La herramienta sirve como prototipo académico; las instrucciones de seguridad deben confirmarse en el manual original.

\section{Chatbot para los documentos del curso}

Se creó un segundo RAG, independiente del industrial, sobre los siete documentos de \texttt{03/data/mds}: introducción a la IA, autómatas celulares, algoritmos genéticos, programación genética, redes neuronales, Machine Learning e IA generativa.

El sistema indexa los documentos, crea embeddings con \texttt{embeddinggemma}, recupera fragmentos relevantes y genera respuestas con \texttt{gemma3:1b}, incluyendo las fuentes consultadas. La implementación y las instrucciones están en \texttt{03/rag/README.md}.

\begin{verbatim}
python3 03/rag/scripts/build_index.py
python3 03/rag/scripts/embed.py embeddinggemma 4
python3 03/rag/scripts/chat.py "¿Qué diferencia hay entre algoritmos genéticos y programación genética?"
\end{verbatim}

La batería de evaluación contiene \textbf{21 preguntas}, tres por documento. El archivo \texttt{03/preguntas\_y\_respuestas.md} se actualiza después de cada respuesta y puede estar parcialmente completado mientras Ollama trabaja. El progreso se consulta con:

\begin{verbatim}
python3 03/rag/scripts/progreso.py --watch
\end{verbatim}

Para reanudar las preguntas pendientes: \texttt{python3 03/rag/scripts/evaluate\_questions.py}. Las respuestas automáticas deben revisarse antes de presentarlas como contenido académico validado.

\section{Agente de mantenimiento}

Se implementó un agente de consola independiente de los dos RAG. Usa SQLite para persistir equipos, repuestos, existencias, órdenes de trabajo e historial de fallas. Permite buscar un repuesto por código, consultar su cantidad y ubicación, crear una orden con falla y equipo, y registrar cambios de estado en el historial. Crear una orden también registra la falla; los eventos posteriores se conservan para mantener trazabilidad.

El código y las instrucciones están en \texttt{04/agente.py} y \texttt{04/README.md}. El comando \texttt{demo} carga datos ficticios para reproducir el ejercicio:

\begin{verbatim}
python3 04/agente.py demo
python3 04/agente.py buscar-repuesto REP-DIS-01
python3 04/agente.py consultar-almacen REP-DIS-01
python3 04/agente.py crear-orden --equipo EQ-PUL-01 --falla "Vibración excesiva" --repuesto REP-DIS-01
python3 04/agente.py actualizar-historial --orden 1 --estado en_proceso --nota "Disco inspeccionado"
\end{verbatim}

Las pruebas de \texttt{04/pruebas\_y\_resultados.md} documentan comandos y respuestas reales. La última ejecución registró \textbf{7/7 comprobaciones aprobadas}, incluidas cinco pruebas automáticas. Se pueden repetir con \texttt{python3 04/evaluar.py}. En una base que ya contenga órdenes, se debe usar el ID devuelto por \texttt{crear-orden} en lugar del \texttt{1} del ejemplo.

\section{Uso de un LLM con ESP32}

En \texttt{05.md} se analizó la viabilidad de usar un microcontrolador ESP32 con un LLM. La propuesta es utilizar el ESP32 para capturar datos y comunicarse por Wi-Fi con un servidor o API que ejecute el modelo. Se describen las limitaciones de memoria, procesamiento, almacenamiento, conexión, latencia, consumo, costos y privacidad. Este punto es un análisis conceptual; no incluye firmware ni montaje físico.

\section{Requisitos para ejecutar los chatbots}

\begin{itemize}
\item Python 3.10 o posterior y Ollama instalados localmente.
\item Modelos \texttt{embeddinggemma} y \texttt{gemma3:1b} disponibles en Ollama.
\item Servicio de Ollama accesible en \texttt{http://localhost:11434}.
\item Los índices SQLite se generan con los comandos de cada punto; no forman parte del código fuente versionado.
\end{itemize}

El agente del punto 4 solo necesita Python 3. Los detalles de uso y regeneración están en los README de cada carpeta.


\end{document}
